\clearpage
\fancyhead[L]{\small BODE. ASTRONOMISCHES JAHRBUCH.}
\hypertarget{native-630-lower}{}
\begin{center}\small Göttingische gelehrte Anzeigen. Stück 5. Seiten 41..48. 8. Januar 1821.\\[3pt]\rule{.72\linewidth}{.3pt}\end{center}

\emph{Astronomisches Jahrbuch für das Jahr \textup{1823}, mit einer Sammlung der neuesten Abhandlungen, Beobachtungen und Nachrichten aus den astronomischen Wissenschaften. Berechnet und herausgegeben von} Dr. J. E. \textsc{Bode}, königlicher Astronom usw. Berlin 1820. Erhältlich beim Verfasser und in Kommission bei F. Dümmler. 252 Seiten im Oktavformat, dazu eine Kupfertafel.

\begingroup\brokenpenalty=10000\relax
Jede Ausgabe des Astronomischen Jahrbuchs ist den Freunden der Astronomie derzeit doppelt willkommen. Seit die von den Herren \textsc{von Lindenau} und \textsc{Bohnenberger} herausgegebene astronomische Zeitschrift ihr Erscheinen eingestellt hat, bietet das Jahrbuch in Deutschland das einzige Forum für diese Wissenschaft. Die diesem Jahrgang beigefügten Aufsätze sind die folgenden. Versuch über die physische Beschaffenheit der Kometen, insbesondere des Kometenschweifs, von Herrn Prof. \textsc{Fischer}. Mit Vergnügen sieht man hier zusammengestellt, was die zuverlässigsten Beobachtungen und geläuterten Auffassungen über diesen Gegenstand gelehrt haben. Gegen einige einzelne Behauptungen lassen sich \SourceLeaf{631}jedoch noch mehrere Einwände erheben. So sind bekanntlich die Meinungen der Astronomen darüber, ob Kometen mit eigenem Licht oder mit reflektiertem Sonnenlicht leuchten, noch geteilt. Ein Astronom, der sich vielfach und in hohem Maße um die Theorie der Kometen verdient gemacht hat, hat in einem früheren Aufsatz so gut wie alles behandelt, was sich nach dem gegenwärtigen Stand unserer Kenntnisse darüber sagen lässt. Der Verfasser vertritt die erste Auffassung, weil das Licht der kleinsten Fixsterne selbst durch die dichteren Teile der Kometenschweife dringt und weil einigen Beobachtungen zufolge Fixsterne sogar durch den Kern hindurch gesehen worden sein sollen. Doch scheint er daraus zu Unrecht zu schließen, dass die Kometenmasse Sonnenlicht überhaupt nicht merklich reflektieren könne und dass das Licht, in dem wir die Kometen sehen, notwendigerweise ganz ihr eigenes sein müsse. Jene Beobachtung beweist nur, dass die Kometenmasse bei Weitem mehr Licht durchlässt als reflektiert. Bedenkt man aber, um wie viel blasser selbst die hellsten Kometen als Planetenscheiben erscheinen, so braucht man selbst dann kein Eigenlicht der Kometen anzunehmen, wenn nur ein äußerst geringer Teil des Sonnenlichts zurückgeworfen wird. Bei seiner Annahme, dass die Kometenmassen Stoffe enthalten, die gegenüber der Sonne negativ gravitieren, scheint dem Verfasser entgangen zu sein, dass die Bewegungen solcher Kometen dann nicht genau den Gesetzen \textsc{Kepler}s folgen würden. Auch wenn wir nicht bestreiten, dass eines der Gesetze \textsc{Kepler}s bei der Bewegung der Kometen tatsächlich eine gewisse Änderung erfahren könnte und dass dieser Umstand eigentlich bei keinem Kometen systematisch untersucht worden ist, lässt sich doch aus allen bisherigen Berechnungen schließen, dass eine solche Änderung nur äußerst klein sein dürfte. --- Beobachtungen des Kometen von 1819 an der Sternwarte Bogenhausen bei München von Herrn Steuerrat \textsc{Soldner}. --- Geografische Ortsbestimmungen in Ostfriesland von Herrn Prof. \textsc{Oltmanns}. --- Über die aus astronomischen Beobachtungen bestimmte Länge von Pisa, von Herrn Prof. \textsc{Wurm}. Das Ergebnis dieser wertvollen Untersuchung ist, dass der \emph{Unterschied} zwischen der von \textsc{Inghirami} aus geodätischen Messungen und der aus den astronomischen Beobachtungen abgeleiteten Länge wegen der schlechten Qualität der letzteren Beobachtungen eigentlich gar nichts beweisen kann. --- Von demselben Astronomen stammen Beiträge zur geografischen Längenbestimmung aus Beobachtungen der Sonnenfinsternisse vom 18. November 1816 und 4. Mai 1818. --- Astronomische Beobachtungen in Wilna von Herrn Prof. \textsc{Sniadecky}, in Palermo von Herrn \textsc{Cacciatore} (die Kometen von 1819 betreffend) und in Prag von den Herren Professoren \textsc{David} und \textsc{Bittner}. --- Noch etwas über den großen Kometen und seinen Vorübergang vor der Sonne von Dr. \textsc{Olbers}. Die schon im astronomischen Jahrbuch für 1822 angeführte Beobachtung des Generals \textsc{von Lindener}, der am 26. Juni 1819 in der Stunde, in der der Komet mit Sicherheit vor der Sonnenscheibe stehen musste, die Sonne ganz ohne Flecken gesehen hatte, schien vollständig zu beweisen, was ohnehin höchst wahrscheinlich war, nämlich dass die lockere Kometenmasse viel zu wenig Sonnenlicht auffangen könnte, um die geringste Trübung zu verursachen. Dennoch zeigt dieser Vorfall auf bemerkenswerte Weise, wie leicht man sich in solchen Fällen irren kann, wenn man nicht besonders aufmerksam beobachtet. Mehrere andere Beobachter haben nämlich am selben Tag tatsächlich Sonnenflecken gesehen, darunter Prof. \textsc{Schumacher}, damals in Altona, und Prof. \textsc{Brandes} in Breslau. Die Beobachtungen des Herrn \textsc{von Lindener} verlieren daher freilich ihre Beweiskraft. Denn wenn ihm die wirklichen Sonnenflecken entgangen sind, konnte er erst recht nicht den Kometen bemerken, selbst wenn dieser mit stärkeren Instrumenten oder bei gesteigerter Aufmerksamkeit erkennbar gewesen wäre. Und tatsächlich könnten zwei andere Beobachtungen beinahe dazu bewegen, Letzteres anzunehmen, wären da nicht einige Begleitumstände, die die Sache wieder ganz ungewiss machten. Dr. \textsc{Gruithuisen} in München sah nämlich an genau diesem Tag und in der Stunde, in der der Komet vor der Sonne stand, außer zwei anderen gewöhnlichen Sonnenflecken noch einen sehr kleinen ohne deutliche Begrenzung, den er an \SourceLeaf{632} früheren Tagen nicht bemerkt hatte. Dr. \textsc{Gruithuisen} beschreibt ihn jedoch als einen \emph{schwarzen} Punkt, und \emph{so} konnte sich dieser Komet doch nicht zeigen. Daher wären wir eher geneigt, die Erscheinung für einen gewöhnlichen Sonnenfleck zu halten, der erst kurz zuvor entstanden war und in den schwächeren Fernrohren der Herren \textsc{Schumacher} und \textsc{Brandes} unsichtbar blieb. Auch Prof. \textsc{Wildt} in Hannover bemerkte um dieselbe Morgenstunde einen sehr verwaschenen Fleck auf der Sonne. Was das Datum betrifft, ist er sich jedoch selbst nicht sicher, da er erst viel später und nur aus dem Gedächtnis davon berichtete. Die Hauptfrage bleibt unter diesen Umständen weiterhin offen und wird es vielleicht noch lange bleiben, da solche Konjunktionen äußerst selten sind. --- Die Ephemeride des Nordsterns für alle oberen Kulminationen des Jahres 1821 wird allen praktischen Astronomen sehr willkommen sein, die die kleine Druckschrift der Herren \textsc{Struve} und \textsc{Walbeck}, aus der sie entnommen ist, nicht selbst besitzen. Eine Seite mehr hätte sie wohl verdient, damit die letzte Dezimalstelle nicht hätte entfallen müssen. --- Astronomische Beobachtungen in Wien von Prof. \textsc{Bürg}, in Prag von Prof. \textsc{Hallaschka} und in Berlin vom Herausgeber. --- Über die beobachtete Existenz einer Photosphäre der Venus im Jahr 1820 von Geheimrat \textsc{Pastorf}. Herr P. sah im April dieses Jahres die Venus (und später auch Jupiter) von einem kreisförmigen, sehr scharf begrenzten Lichtschimmer umgeben, den er als eine zum Planeten selbst gehörende Lichtsphäre ansieht. Diese Erscheinung verdient eine genauere Untersuchung, weil bei solchen Gegenständen eine optische Täuschung so leicht eine Rolle spielen kann. Herr \textsc{Pastorf} stellte im April fest, dass der Durchmesser dieses Lichtschimmers, der ähnlich aussieht wie die Nachtseite des Mondes im Erdlicht, 16 Bogenminuten betrug. Bei der von Herrn \textsc{Pastorf} am 27. April beobachteten Erscheinung trat ein kleiner teleskopischer Stern oben am \emph{östlichen} Rand der Lichtsphäre der Venus ein, blieb \emph{einige Minuten} unsichtbar und erschien dann an der \emph{Westseite} wieder. Etwas bedenklich ist dabei, dass die Venus an diesem Tag in jeder Minute nur um 3 Bogensekunden vorrückte. Herr \textsc{Pastorf} stellt die Hypothese auf, dass die Venus einen Trabanten haben könnte, der nicht über die Lichtsphäre hinauskommt und der, wie Herr \textsc{Pastorf} glaubt, bei einem Durchmesser von nur 3 Bogensekunden für uns im reflektierten Sonnenlicht immer unsichtbar bleiben müsste, außer wenn er zuweilen wie ein dunkler Fleck vor der Venus erschiene. Wir wollen dem Urteil unserer Leser darüber nicht vorgreifen. --- Beobachtete Rektaszensionen von Saturn und Vesta im Jahr 1819 sowie von Pallas und Mars im Jahr 1820 am Meridianfernrohr der Göttinger Sternwarte, von Hofrat \textsc{Gauss}. --- Beschreibung des an der Königsberger Sternwarte aufgestellten \textsc{Reichenbach}'schen Meridiankreises von Prof. \textsc{Bessel}. --- Beobachtungen des Kometen von 1819 nebst Sternbedeckungen von Prof. \textsc{Struve}. --- Bestimmung der Schiefe der Ekliptik mit einem \textsc{Reichenbach}'schen Meridiankreis an der Sternwarte Bogenhausen bei München von Herrn \textsc{Soldner}. Die dafür verwendeten Beobachtungen liefern zugleich den Sonnenhalbmesser, und zwar $16'\,0''93$ für die mittlere Entfernung. Wir bemerken bei dieser Gelegenheit, dass sich aus 271 von Januar bis Juli 1820 am \textsc{Reichenbach}'schen Mittagsfernrohr der hiesigen Sternwarte beobachteten Fadenantritten ein fast genau ebenso großer Halbmesser ergeben hat, nämlich $16'\,1''01$. --- Astronomische Beobachtungen in Kremsmünster von Herrn \textsc{Derflinger}. --- Beobachtungen von Juno, Pallas, Ceres, Mars und Uranus sowie der Schiefe der Ekliptik von Herrn Prof. \textsc{Nicolai} in Mannheim. Letztere, mit dem dreifüßigen \textsc{Reichenbach}'schen Repetitionskreis bestimmt, ergibt sich aus der Wintersonnenwende um $6''$ kleiner als aus der Sommersonnenwende, wodurch, ebenso wie durch die mit demselben Instrument bestimmten Sterndeklinationen, erneut bestätigt wird, dass man jedes astronomische Instrument bei der Betrachtung der Einwirkung der Schwerkraft auf seine Teile als Individuum behandeln muss. --- Über die Genauigkeit der Beobachtungen am Mittagsfernrohr von Herrn Dr. \textsc{Walbeck}. Hier wird zum ersten Mal die richtige Behandlung dieses \SourceLeaf{633} wichtigen Gegenstands gelehrt und auf die hervorragenden Dorpater Beobachtungen angewandt. Auf S.\,186 Z.\,3 von unten steht durch einen Druckfehler tang $\delta$ statt sec $\delta$. --- Astronomische Beobachtungen in Hamburg von Herrn \textsc{Rümker}. --- Über die geografische Lage von Dresden von Herrn Dr. \textsc{Raschig}. --- Astronomische Bemerkungen von Herrn Prediger \textsc{Luthmer} in Hannover. --- Über das wahre Datum der nächtlichen Schlacht am Halys von Herrn Prof. \textsc{Oltmanns}. In dieser hervorragenden Abhandlung, bei der wir ungern sehen, dass aus Platzgründen alle Zitate und Anmerkungen weggelassen wurden, wird sehr überzeugend gezeigt, dass die viel besprochene Sonnenfinsternis während der berühmten Schlacht nur die vom 30. September 609 v. Chr. gewesen sein kann. --- Über die Bahn des Pons'schen (\textsc{Enke}'schen) Kometen samt Berechnung seines Laufs bei seiner nächsten Wiederkehr im Jahr 1822 von Herrn Prof. \textsc{Enke}. Die wiederholte und vollständige geführte Berechnung der Störungen, die dieser Komet von 1786--1819 erfahren hat, gibt die Beobachtungen bei den vier bisherigen Erscheinungen äußerst gut wieder. Allerdings zeigt sich eine bemerkenswerte Beschleunigung der Bewegung, die sich durch die Rechnung noch nicht erklären lässt. Die drei Umläufe von 1786--1795 ergeben eine um einen halben Tag längere und die vier Umläufe von 1805--1819 eine um einen halben Tag kürzere Umlaufzeit als die drei von 1795--1805. Sollte die nächste Wiedererscheinung des Kometen ein ähnliches Ergebnis liefern, wird man über die Ursache dieses Phänomens plausible Vermutungen äußern können. Herr Prof. \textsc{Enke} hat mit einer praktischen Ephemeride bestens dafür gesorgt, dass der Komet im Jahr 1822, in dem er am 24. oder 25. Mai durch sein Perihel gehen wird, leichter wiedergefunden werden kann. Auf der Südhalbkugel wird die Beobachtung im Juni und Juli keine Schwierigkeiten bereiten. In Europa wird er jedoch um diese Zeit wegen seiner südlichen Lage überhaupt nicht zu sehen sein. Früher wird es wegen seiner großen Lichtschwäche äußerst schwierig, wenn nicht unmöglich sein, ihn zu erkennen. Damit aber nichts unversucht bleibt, wünschen wir, dass der um diesen Kometen so verdiente Astronom die am 25. Februar beginnende Ephemeride noch um ein paar frühere Monate ergänzt. Wegen des hohen Stands des Kometen bei dunkler Nacht besteht früher trotz seiner noch größeren Lichtschwäche vielleicht etwas mehr Hoffnung als zu der Zeit, in der er nach Ende der Dämmerung schon so nahe am Horizont steht. --- Auch für den kleineren Kometen des Jahres 1819 hat Herr Prof. \textsc{Enke} eine elliptische Bahn mit einer Umlaufzeit von nur $5\tfrac{1}{2}$ Jahren gefunden, die sowohl die Marseiller als auch die erst später bekannt gewordenen Mailänder Beobachtungen hervorragend wiedergibt. --- Vom selben Astronomen, die Opposition der Vesta von 1819 und eine Ephemeride für die nächste Erscheinung dieses Planeten. --- Ephemeriden für Juno und Pallas 1821 von Herrn Prof. \textsc{Nicolai} und von Herrn \textsc{von Staudt} in Göttingen. --- Astronomische Beobachtungen im Jahr 1820 von Herrn Hofrat \textsc{Gauss}. --- Über die Bestimmung der geografischen Breite mithilfe des Polarsterns von Herrn Prof. \textsc{Dirksen} in Berlin (Analyse einer praktischen Tafel, die Herr Prof. \textsc{Schumacher} zu diesem Zweck angegeben hat). --- Weitere Beobachtungen von Sternbedeckungen, Verfinsterungen der Jupitermonde und der großen Sonnenfinsternis vom 7. September 1820 durch Herrn Prof. \textsc{Rümker} in Hamburg sowie Beobachtungen der letzteren durch Herrn Prof. \textsc{Nicolai} und Herrn \textsc{von Heiligenstein} in Mannheim. --- Von den kurzen Nachrichten am Ende dieses Bandes heben wir noch die über die in Abo errichtete neue Sternwarte und die über die in den russischen Ostseeprovinzen vorzunehmende Gradmessung hervor. --- Die Kupfertafel zeigt außer den Sternbedeckungen und den beiden Mondfinsternissen des Jahres 1823 den berechneten Vorübergang des großen Kometen von 1819 vor der Sonne, die beobachtete geozentrische Bewegung des \textsc{Enke}'schen Kometen 1818 und 1819 und die Venus mit einem von Herrn \textsc{Pastorff} beobachteten Fleck.
\par\clearpage\endgroup
